\chapterdivider{4}{Public API and JAX execution model}
  {How does a scientific model become a transformable branch computation?}
  {Featured application: one functional API used eagerly, under JIT, and across batches}

\begin{frame}[fragile]{The front door is one functional call}
\releasedbadge
\begin{lstlisting}
problem = jc.bif_problem(rhs, u0, p0, args=args)
result = jc.continuation(
    problem,
    jc.PseudoArclength(),
    p_span=(p_min, p_max),
    settings=jc.ContinuationPar(ds=0.02, max_steps=300),
    events=[jc.Fold(), jc.Hopf()],
)
\end{lstlisting}
\begin{block}{The call contract}
The problem holds the model and starting state; the algorithm chooses how to
follow the branch; settings bound the computation; events request diagnostics.
The return value is always a \texttt{ContinuationResult}.
\end{block}
\begin{alertblock}{Starting pair}
Continuation begins from \texttt{(problem.u0, p\_span[0])}; supply a state that
solves the model at that starting parameter.
\end{alertblock}
\end{frame}

\begin{frame}{Read the scientific result as a small tree}
\centering
\begin{tikzpicture}[node distance=5mm and 10mm,>=Latex,
 root/.style={rounded corners,draw=JaxBlue,very thick,fill=SoftBlue,
              minimum width=34mm,minimum height=11mm,align=center},
 group/.style={rounded corners,draw=JaxGreen,thick,fill=JaxGreen!10,
               minimum width=30mm,minimum height=10mm,align=center},
 leaf/.style={rounded corners,draw=JaxOrange,thick,fill=JaxOrange!10,
              minimum width=46mm,minimum height=10mm,align=left}]
\node[root] (result) {\texttt{ContinuationResult}};
\node[group,right=15mm of result,yshift=14mm] (branch) {\texttt{result.branch}};
\node[group,right=15mm of result,yshift=-14mm] (events) {\texttt{result.events}};
\node[leaf,right=of branch,yshift=11mm] (arrays)
  {\texttt{.states / .params}\\\texttt{.stable / .eigenvalues}};
\node[leaf,right=of branch,yshift=-11mm] (valid)
  {\texttt{.valid / .n\_valid}};
\node[leaf,right=of events] (hit)
  {each hit: \texttt{.kind / .u / .p / .info}};
\draw[->,thick] (result)--(branch);
\draw[->,thick] (result)--(events);
\draw[->,thick] (branch)--(arrays);
\draw[->,thick] (branch)--(valid);
\draw[->,thick] (events)--(hit);
\end{tikzpicture}
\vspace{3mm}
\begin{block}{Translate names once}
\texttt{stable} is the stability flag; event fields \texttt{u}, \texttt{p},
and \texttt{info} are the event state, parameter, and diagnostics.
\end{block}
\end{frame}

\begin{frame}{One result has an eager view and a traced view}
\begin{columns}[T]
\column{.48\textwidth}
\begin{block}{Eager call: convenient arrays}
Branch arrays are trimmed to the computed points, \texttt{branch.valid} is
\texttt{None}, and \texttt{branch.n\_valid} is the trimmed length. Event hits
may be returned as Python objects.
\end{block}
\column{.48\textwidth}
\begin{block}{Inside \texttt{jit} / \texttt{vmap}: fixed arrays}
Branch arrays retain all \texttt{max\_steps + 1} slots.
\texttt{branch.valid} selects real points, while
\texttt{result.stats["n\_valid"]} records their count.
\end{block}
\end{columns}
\vspace{4mm}
\begin{center}
\begin{tikzpicture}[x=.72cm,y=.62cm]
\foreach \i in {0,...,12} {\draw[fill=SoftGray] (\i,0) rectangle +(0.68,0.72);}
\foreach \i in {0,...,7} {\draw[fill=JaxBlue!45] (\i,0) rectangle +(0.68,0.72);}
\draw[JaxGreen,very thick] (-.05,-.12) rectangle (5.72,.84);
\node[below] at (2.85,-.15) {\texttt{valid = True}};
\node[below] at (7.3,-.15) {reserved capacity};
\end{tikzpicture}
\end{center}
\begin{alertblock}{Current transformation boundary}
Use \texttt{events=()} inside \texttt{jit}/\texttt{vmap}; event orchestration is
eager-only even though branch and stability arrays are transformation-safe.
\end{alertblock}
\end{frame}

\begin{frame}{One public call selects one of two compiled engines}
\centering
\begin{tikzpicture}[node distance=8mm and 9mm,>=Latex,
 box/.style={rounded corners,draw=JaxBlue,thick,fill=SoftBlue,
             minimum width=31mm,minimum height=11mm,align=center},
 engine/.style={rounded corners,draw=JaxGreen,very thick,fill=JaxGreen!10,
                minimum width=70mm,minimum height=15mm,align=left}]
\node[box] (model) {\texttt{rhs, u0, p0, args}};
\node[box,right=of model] (problem) {\texttt{BifProblem}};
\node[box,right=of problem] (call) {\texttt{continuation(...)}};
\node[engine,below=of call] (backends)
  {\texttt{Natural()} $\to$ \texttt{natural\_scan}\\
   \texttt{PseudoArclength()} $\to$ \texttt{pseudo\_arclength\_scan}};
\node[box,left=of backends] (result) {\texttt{ContinuationResult}};
\draw[->,thick] (model)--(problem);
\draw[->,thick] (problem)--(call);
\draw[->,thick] (call)--(backends);
\draw[->,thick] (backends)--(result);
\end{tikzpicture}
\vspace{4mm}
\begin{block}{Why the boundary matters}
Users compose immutable problem, algorithm, settings, and event values.
\texttt{continuation} dispatches the algorithm to one of two compiled engines;
both return the same result contract to demos and downstream analysis.
\end{block}
\vfill
{\scriptsize Carry fields, array-selection mechanics, and reassembly details are implementation-appendix material (Chapter 12).}
\end{frame}

\begin{frame}{``Scan'' names the engine, not the loop primitive}
\begin{columns}[T]
\column{.47\textwidth}
\begin{block}{Pseudo-arclength}
The numerical method: tangent prediction followed by augmented Newton
correction, so an ordinary fold can be passed.
\end{block}
\column{.47\textwidth}
\begin{block}{Scan engine}
The implementation style: package the bounded branch sweep into one compiled,
functional array program.
\end{block}
\end{columns}
\vspace{5mm}
\begin{alertblock}{Precise statement}
\texttt{pseudo\_arclength\_scan} uses \texttt{jax.lax.while\_loop} for the
outer continuation loop; it does not implement the branch sweep with
\texttt{jax.lax.scan}.
\end{alertblock}
\begin{block}{Nested work is bounded too}
Each continuation attempt calls a Newton loop with a finite iteration budget;
the outer loop has a fixed maximum number of accepted steps.
\end{block}
\end{frame}

\begin{frame}{Whole-loop JIT changes orchestration, not mathematics}
\centering
\begin{tikzpicture}[node distance=6mm and 7mm,>=Latex,
 stage/.style={draw=JaxBlue,fill=SoftBlue,rounded corners,minimum width=24mm,
              minimum height=9mm,align=center},
 boundary/.style={draw=JaxGreen,very thick,rounded corners,inner sep=5mm}]
\node[stage] (tan) {compute\\tangent};
\node[stage,right=of tan] (pred) {predict};
\node[stage,right=of pred] (corr) {Newton\\correct};
\node[stage,right=of corr] (adapt) {accept, store,\\adapt};
\draw[->,thick] (tan)--(pred);
\draw[->,thick] (pred)--(corr);
\draw[->,thick] (corr)--(adapt);
\draw[->,thick] (adapt.south) -- ++(0,-.55) -| (tan.south);
\node[boundary,fit=(tan)(pred)(corr)(adapt),
      label={[text=JaxGreen!60!black,font=\bfseries]above:one compiled continuation program}] {};
\end{tikzpicture}
\vspace{4mm}
\begin{columns}[T]
\column{.48\textwidth}
\begin{alertblock}{Cold call}
Tracing and compilation are included. Report this cost separately and wait for
device completion before timing.
\end{alertblock}
\column{.48\textwidth}
\begin{block}{Warm call}
A compatible call reuses the executable and measures steady-state execution.
\end{block}
\end{columns}
\vspace{2mm}
{\small The model function, \texttt{max\_steps}, and linear solver are static;
changing them or array shapes/dtypes may trigger recompilation.}
\end{frame}

\begin{frame}[fragile]{\texttt{vmap} batches independent branches}
\begin{lstlisting}
def run_one(design_value):
    problem = jc.bif_problem(
        rhs, u0, p0, args=design_value
    )
    return jc.continuation(
        problem, jc.PseudoArclength(),
        p_span=(p_min, p_max), settings=settings,
        events=(),
    )

branches = jax.vmap(run_one)(design_values)
\end{lstlisting}
\begin{block}{What is parallel}
Each batch element owns one sequential continuation history. Fixed buffers give
every branch the same outer shape; \texttt{branches.branch.valid} identifies
the actual points in each history.
\end{block}
\begin{alertblock}{What is not parallel}
\texttt{vmap} does not parallelize the predictor--corrector steps within one
branch, because step $k+1$ depends on the accepted point at step $k$.
\end{alertblock}
\end{frame}

\begin{frame}{Autodiff serves local equations and converged roots}
\begin{columns}[T]
\column{.48\textwidth}
\begin{block}{During branch computation}
Automatic Jacobians provide $F_{\state}$ and $F_p$ for tangents, Newton
systems, stability, and forward sensitivities through the bounded sweep.
\end{block}
\column{.48\textwidth}
\begin{block}{At a selected implicit root}
A custom VJP differentiates the converged equation $G(x^*,\theta)=0$ without
reverse-mode backpropagation through Newton iterations:
\[
  \frac{dx^*}{d\theta}=-G_x^{-1}G_\theta.
\]
\end{block}
\end{columns}
\vspace{4mm}
\begin{alertblock}{Derivatives inherit numerical quality}
Autodiff differentiates the implemented equations. It does not repair poor
conditioning, select the intended branch, or make a failed root solve converge.
\end{alertblock}
\end{frame}

\begin{frame}{Solver protocols isolate numerical policy}
\centering
\begin{tikzpicture}[node distance=12mm and 18mm,>=Latex,
 box/.style={rounded corners,draw=JaxBlue,thick,fill=SoftBlue,
             minimum width=40mm,minimum height=12mm,align=center},
 seam/.style={rounded corners,draw=JaxGreen,thick,fill=JaxGreen!10,
              minimum width=45mm,minimum height=12mm,align=center}]
\node[box] (engine) {continuation correction};
\node[box,below=of engine] (stable) {stability post-pass};
\node[seam,right=of engine] (linear)
  {\texttt{LinearSolver(A, b)} $\to x$\\shipped default: \texttt{Dense()}};
\node[seam,right=of stable] (eigen)
  {\texttt{EigenSolver(A)} $\to \lambda$\\shipped default: \texttt{DenseEigen()}};
\draw[->,thick] (engine)--(linear);
\draw[->,thick] (stable)--(eigen);
\end{tikzpicture}
\vspace{6mm}
\begin{block}{Architectural payoff}
The branch engine depends on small callable contracts rather than hard-coding
linear solves or eigensolvers. The public \texttt{jc.Solvers} bundle supplies
the two strategies to \texttt{jc.continuation}.
\end{block}
\end{frame}

\begin{frame}{Bounded termination is part of the execution contract}
\begin{columns}[T]
\column{.48\textwidth}
\begin{block}{The outer branch loop stops when}
\begin{itemize}
  \item the requested parameter endpoint is reached;
  \item the accepted-step budget is exhausted;
  \item a failed correction reaches \texttt{ds\_min};
  \item the accepted state becomes non-finite.
\end{itemize}
\end{block}
\column{.48\textwidth}
\begin{block}{Each Newton loop stops when}
\begin{itemize}
  \item the residual meets tolerance;
  \item a residual becomes non-finite;
  \item \texttt{newton\_max\_iter} is reached.
\end{itemize}
\end{block}
\end{columns}
\vspace{4mm}
\begin{alertblock}{Bounded does not mean successful}
Termination prevents an unbounded compiled computation. Scientific success
still requires inspecting validity, residuals, events, and branch coverage.
\end{alertblock}
\end{frame}

\begin{frame}{JAX mechanisms enable capabilities, not guarantees}
\scriptsize
\begin{center}
\begin{tabular}{>{\raggedright\arraybackslash}p{.20\textwidth}
                >{\raggedright\arraybackslash}p{.31\textwidth}
                >{\raggedright\arraybackslash}p{.38\textwidth}}
\toprule
JAX mechanism & Enables & Does not guarantee\\
\midrule
autodiff & Jacobians and sensitivities & good conditioning\\
JIT + \texttt{while\_loop} & compiled branch execution & a more accurate method\\
fixed shapes & transformation compatibility & every event path is batchable\\
\texttt{vmap} & independent ensemble branches & parallel steps within one branch\\
custom VJP root & implicit reverse-mode gradients & convergence from a poor seed\\
\bottomrule
\end{tabular}
\end{center}
\vspace{4mm}
\begin{block}{The practical reading}
JAX removes derivative bookkeeping and Python/device orchestration from a
well-posed computation. Numerical method choice, convergence, conditioning,
and validation remain scientific responsibilities.
\end{block}
\end{frame}
